\scanpage{21}
\mspage{9}
% Folio d'archive : 20. Feuille double ; lecture de la moitié gauche puis droite.
\uline{Ex 4} Autre cas particulier. Soit
\[
M=\text{complexes d'une catégorie abélienne},\ \mathcal A
\qquad W\ \text{les quasi-isom.}
\]
\[
\uline{\operatorname{Hom}}(A,M)=\text{complexes de }\uline{\operatorname{Hom}}(A,\mathcal A)
\]
\[
W_A=\text{quasi-isom d'iceux.}
\]
\inserted{$\mathcal H(A)$ catégorie dérivée de $\uline{\operatorname{Hom}}(A,\mathcal A)$}

Si p. ex. $\mathcal A=$ faisceaux de $\unsure{k|h}$-modules sur topos
\deleted{$\unsure{\mathcal E|\mathcal T}$} $\mathcal T$
alors
\[
\uline{\operatorname{Hom}}(A,\mathcal A)=
\text{faisceaux de $\unsure{k|h}$-modules sur }\inserted{\text{le topos }}
\widehat{A^\circ}\mathbin{\times_{\mathrm{top}}}\mathcal T.
\]
Et $f^*:\uline{\operatorname{Hom}}(B,M)\longrightarrow
\uline{\operatorname{Hom}}(A,M)$ est le foncteur image inverse associé à
\diagram{\begin{tikzcd}
\widehat{A^\circ}\times\mathcal T
\arrow[d,"{f^\circ_{\mathrm{top}}\times\operatorname{id}_{\mathcal T}}"]\\
\widehat{B^\circ}\times\mathcal T
\end{tikzcd}}
\margin{induisant $\unsure{(L f^{\circ *})|(L f^{\circ *}_{\mathrm{top}})}$}
Le foncteur \inserted{induit sur les catégories dérivées}
a un adjoint à \uline{droite} $R f_*$,
\deleted{mais} généralement pas d'adjoint à gauche
\inserted{sans doute, $=$ si $\mathcal T$ : topos ponctuel ou
$\mathcal A=(\unsure{k|h}\text{-modules})$}.
Si pourtant, car au niveau des \deleted{topos} faisceaux de modules
$(f^\circ_{\mathrm{top}})_!$ existe et est exact !
\unsure{pas dur|pas abs.}, donc
% Suite, en haut de la moitié droite.
au niveau des \inserted{complexes de} faisceaux abéliens il
\deleted{passe}\inserted{transforme} quasi-isom. en iso, donc
passe aux catégories dérivées en $L f_!$.

\figure{\begin{tikzpicture}\draw (0,0)--(2,0);\end{tikzpicture}}
\uline{Propriétés formelles} supplémentaires
\deleted{\unsure{de la théorie}}
\margin{\unsure{les deux propriétés impliquent la commutation aux sommes et produits des $g_{\mathcal H}^*$}}

% Hot 4 entouré.
\figure{\begin{tikzpicture}\node[draw,ellipse] {Hot 4};\end{tikzpicture}}
\diagram{\begin{tikzcd}[column sep=large,row sep=large]
A\arrow[d,"f"'] & A'\arrow[l,"g'"']\arrow[d,"f'"]\\
B & B'\arrow[l,"g"]
\arrow[ul,phantom,"\mathrm{cart.}"]
\end{tikzcd}}
\[
g^* f_!\xleftarrow{\sim}f'_!g'^*
\quad\text{si }f\text{ cofibrant}\quad[\text{ou }g\text{ fibrant ??}]
\]
\[
g^* f_*\xrightarrow{\sim}f'_*g'^*
\quad\text{si }f\text{ fibrant}\quad[\text{ou }g\text{ cofibrant ??}]
\]
Les formules \inserted{sont} ``transposées'' \inserted{l'une de l'autre},
par passage aux adjoints (à droite, resp. à g.)

\figure{\begin{tikzpicture}\node[draw,ellipse] {Hot 5};\end{tikzpicture}}
Si $A\xrightarrow{f}B$ pl. fid., alors
\[
f_{\mathcal H}^*:\mathcal H(B)\longrightarrow\mathcal H(A)
\quad\text{foncteur localisation}
\]
i.e.
\[
f_*^{\mathcal H},f_!^{\mathcal H}:\mathcal H(A)\longrightarrow\mathcal H(B)
\quad\text{sont pl. fid.}
\]
(et $f^*f_!\simeq\operatorname{id}$, $f^*f_*\simeq\operatorname{id}$)

\figure{\begin{tikzpicture}\node[draw,ellipse] {Hot 6};\end{tikzpicture}}
Si $f:A\longrightarrow B$ localisation, alors
\[
f_{\mathcal H}^*:\mathcal H(B)\longrightarrow\mathcal H(A)
\quad\text{pl. fid., i.e.}
\]
$f_!^{\mathcal H},f_*^{\mathcal H}$ foncteurs loc. \deleted{iso} et
\[
f_!f^*\simeq f_*f^*\simeq\operatorname{id}
\]
